\documentclass[12pt,a4paper]{article}

\usepackage[spanish]{babel}
\usepackage[utf8]{inputenc}
\usepackage[T1]{fontenc}
\usepackage[top=2.5cm,bottom=2.5cm,left=3cm,right=3cm]{geometry}
\usepackage{parskip}
\usepackage{titlesec}
\usepackage{enumitem}
\usepackage{booktabs}
\usepackage{array}
\usepackage{longtable}
\usepackage{xcolor}
\usepackage{listings}
\usepackage{tikz}
\usetikzlibrary{positioning,shapes.geometric,arrows.meta,calc}
\usepackage[hidelinks]{hyperref}
\usepackage{fancyhdr}

% ---- Colores ----
\definecolor{clssblue}{RGB}{207,226,255}
\definecolor{ifcyellow}{RGB}{255,250,200}
\definecolor{codegreen}{rgb}{0,0.6,0}
\definecolor{backcolour}{rgb}{0.97,0.97,0.97}
\definecolor{codekw}{rgb}{0.13,0.13,0.68}

% ---- Estilos TikZ ----
\tikzset{
  mybox/.style={
    rectangle,
    draw=black!60,
    minimum width=2.8cm,
    minimum height=0.85cm,
    align=center,
    font=\small\sffamily,
    rounded corners=3pt,
  },
  classbox/.style={mybox, fill=clssblue},
  ifcbox/.style={mybox, fill=ifcyellow},
  uses/.style={->, >=Stealth, thick},
  impl/.style={->, >=Stealth, dashed, thick},
  flow/.style={->, >=Stealth, thick},
}

% ---- Estilo código ----
\lstdefinestyle{gostyle}{
  backgroundcolor=\color{backcolour},
  commentstyle=\color{codegreen}\itshape,
  keywordstyle=\color{codekw}\bfseries,
  basicstyle=\ttfamily\footnotesize,
  breaklines=true,
  frame=single,
  rulecolor=\color{black!25},
  numbers=left,
  numberstyle=\tiny\color{gray},
  numbersep=6pt,
  tabsize=2,
  language=Go,
  morekeywords={string,error,bool,int,uuid,context},
}

% ---- Secciones ----
\titleformat{\section}{\large\bfseries\sffamily}{\thesection.}{0.5em}{}[\titlerule]
\titleformat{\subsection}{\normalsize\bfseries\sffamily}{\thesubsection.}{0.5em}{}
\titleformat{\subsubsection}{\small\bfseries\sffamily}{\thesubsubsection.}{0.4em}{}

% ---- Encabezado/pie ----
\pagestyle{fancy}
\fancyhf{}
\fancyhead[L]{\small\sffamily AutoMatch}
\fancyhead[R]{\small\sffamily Diseño de Software --- UCT}
\fancyfoot[C]{\small\thepage}
\renewcommand{\headrulewidth}{0.4pt}

\setlist[itemize]{noitemsep,topsep=3pt,leftmargin=*}
\setlist[enumerate]{noitemsep,topsep=3pt,leftmargin=*}

\begin{document}

% ===== PORTADA =====
\begin{titlepage}
  \thispagestyle{empty}
  \centering
  \vspace*{2.5cm}
  {\Huge\bfseries\sffamily AutoMatch\par}
  \vspace{0.5cm}
  {\large\sffamily Plataforma de Matchmaking para Vehículos Usados\par}
  \vspace{1.2cm}
  \rule{\textwidth}{0.4pt}
  \vspace{0.4cm}
  {\normalsize\sffamily Informe de Proyecto --- Diseño de Software\par}
  {\normalsize\sffamily Universidad Católica de Temuco\par}
  \vspace{0.4cm}
  \rule{\textwidth}{0.4pt}
  \vspace{2cm}
  \begin{tabular}{ll}
    \toprule
    \multicolumn{2}{c}{\textbf{\sffamily Integrantes}} \\
    \midrule
    Lizardo Salazar            & \texttt{lsalazar2021@alu.uct.cl}  \\
    Juan Carrera               & \texttt{juan-xp} (GitHub)          \\
    Benjamín Aliaga Mardones   & \texttt{baliaga2025@alu.uct.cl}    \\
    Benjamín De La Fuente      & \texttt{bdelafuenteh@gmail.com}    \\
    \bottomrule
  \end{tabular}
  \vfill
  {\normalsize\sffamily 2025\par}
\end{titlepage}

\tableofcontents
\newpage

% ===== 1. NOMBRE E INTEGRANTES =====
\section{Nombre e integrantes}

\textbf{Proyecto:} AutoMatch \quad
\textbf{Asignatura:} Diseño de Software \quad
\textbf{Universidad:} Universidad Católica de Temuco

\begin{center}
\begin{tabular}{lll}
  \toprule
  \textbf{Nombre} & \textbf{Usuario Git} & \textbf{Correo} \\
  \midrule
  Lizardo Salazar            & LizardoFSA             & lsalazar2021@alu.uct.cl \\
  Juan Carrera               & juan-xp                & ---                     \\
  Benjamín Aliaga Mardones   & BenjaminAliagaMardones & baliaga2025@alu.uct.cl  \\
  Benjamín De La Fuente      & BenDLF                 & bdelafuenteh@gmail.com  \\
  \bottomrule
\end{tabular}
\end{center}

% ===== 2. PROBLEMA =====
\section{Problema}

En el mercado de vehículos usados, tanto compradores como vendedores enfrentan una fricci\'{o}n
importante al intentar conectarse. Los compradores deben revisar manualmente un gran volumen de
publicaciones que no se ajustan a su presupuesto ni al tipo de vehículo que buscan, invirtiendo
tiempo en opciones irrelevantes. Por su parte, los vendedores publican sus vehículos en
plataformas genéricas donde es difícil distinguir a un comprador genuinamente interesado de uno
que solo navega sin intención real de compra.

\textbf{AutoMatch} propone resolver este problema mediante una plataforma de \emph{matchmaking}
especializada en vehículos. La solución se sustenta en dos pilares:

\begin{enumerate}
  \item \textbf{Personalización del feed}: antes de navegar, el comprador configura su perfil con
        el tipo de vehículo que le interesa (sedán, SUV, hatchback, pickup, van u otro) y un rango
        de presupuesto (mínimo y máximo en CLP). El sistema usa estas preferencias para construir
        un feed filtrado, de forma que el comprador solo ve vehículos que se ajustan a sus
        criterios.
  \item \textbf{Mecanismo de swipe y match automático}: el comprador evalúa cada vehículo del
        feed con un \emph{like} o un \emph{pass}. Cuando hace \emph{like} en un listing activo, el
        sistema crea automáticamente un \emph{match} entre el comprador y el vendedor de ese
        vehículo, y habilita un canal de chat privado entre ellos. Así, la comunicación ocurre
        solo cuando existe interés real de ambas partes, eliminando el ruido para el vendedor y
        agilizando la búsqueda para el comprador.
\end{enumerate}

La plataforma maneja dos roles claramente diferenciados: el rol \textbf{comprador} (\texttt{buyer})
accede al feed y realiza swipes, mientras que el rol \textbf{vendedor} (\texttt{seller}) gestiona
sus publicaciones de vehículos. Cada usuario elige su rol al momento de registrarse.

% ===== 3. AVANCE DEL PROYECTO =====
\section{Avance del proyecto}

El proyecto implementa una aplicación web completa organizada en dos aplicaciones dentro de un
monorepo: \texttt{apps/server} (backend en Go) y \texttt{apps/client} (frontend en Next.js).
Ambas se orquestan junto a la base de datos mediante Docker Compose.

\subsection{Backend --- \texttt{apps/server}}

El backend sigue una \textbf{Clean Architecture} dividida en cuatro capas bien separadas:

\begin{itemize}
  \item \textbf{Dominio} (\texttt{internal/domain}): entidades y lógica de negocio pura, sin
        dependencias de infraestructura. Contiene las entidades \texttt{User}, \texttt{Listing},
        \texttt{Photo}, \texttt{Match}, \texttt{Message}, \texttt{BuyerProfile} y \texttt{Swipe},
        junto con sus constructores validadores y los tipos enumerados
        (\texttt{ListingStatus}, \texttt{Role}, \texttt{SwipeDirection}).
  \item \textbf{Repositorios} (\texttt{internal/repository}): interfaces de acceso a datos e
        implementaciones PostgreSQL. Cada agregado de dominio tiene su propia interfaz
        (\texttt{UserRepository}, \texttt{ListingRepository}, \texttt{MatchRepository},
        \texttt{MessageRepository}, \texttt{ProfileRepository}, \texttt{SwipeRepository}).
  \item \textbf{Servicios} (\texttt{internal/service}): lógica de aplicación que orquesta el
        dominio y los repositorios. Incluye \texttt{AuthService}, \texttt{ListingService},
        \texttt{FeedService}, \texttt{MatchService} y \texttt{ProfileService}.
  \item \textbf{Handlers} (\texttt{internal/handler}): adaptadores HTTP (Gin) que traducen
        peticiones/respuestas y delegan al servicio correspondiente. Incluye
        \texttt{AuthHandler}, \texttt{ListingHandler}, \texttt{FeedHandler},
        \texttt{MatchHandler} y \texttt{ProfileHandler}.
\end{itemize}

Las funcionalidades implementadas son:
\begin{itemize}
  \item \textbf{Autenticación}: registro con validación de email y contraseña (mínimo 6
        caracteres), login con emisión de JWT (HS256, 24 horas de expiración).
  \item \textbf{Gestión de perfiles}: un comprador puede configurar y actualizar su tipo de
        vehículo preferido y rango de presupuesto. Si el perfil no existe, se devuelve uno vacío
        (sin error).
  \item \textbf{Listings (CRUD)}: los vendedores pueden crear publicaciones con marca, modelo,
        año, precio, tipo de vehículo, descripción y hasta 5 fotos; editar campos parcialmente
        (PATCH); y eliminar sus publicaciones. Las fotos se almacenan como URLs en la tabla
        \texttt{listing\_photos}, cargadas en lote para evitar el problema N+1.
  \item \textbf{Feed}: construido dinámicamente por \texttt{FeedService}, que filtra listings
        activos, excluye los ya swipeados por el comprador, excluye los publicados por él mismo y
        aplica las preferencias de su perfil.
  \item \textbf{Swipe y match}: al registrar un swipe de tipo \emph{like}, \texttt{MatchService}
        crea automáticamente un \texttt{Match} en la base de datos. Los swipes de tipo
        \emph{pass} solo se registran sin generar match.
  \item \textbf{Chat}: los participantes de un match (comprador y vendedor del listing asociado)
        pueden enviar y listar mensajes. La autoriz\'{a}ción se verifica a nivel de servicio.
  \item \textbf{CORS}: configurado para aceptar peticiones desde \texttt{localhost:3000} y
        \texttt{localhost:3001} con los métodos GET, POST, PUT, PATCH, DELETE y OPTIONS.
  \item \textbf{Swagger}: documentación interactiva generada con \texttt{swaggo/gin-swagger},
        disponible en \texttt{/swagger/index.html}.
\end{itemize}

\subsection{Frontend --- \texttt{apps/client}}

El cliente es una aplicación Next.js 16 con React 19, TypeScript y Tailwind CSS. La autenticación
se persiste en \texttt{localStorage} y se distribuye en la aplicación mediante un contexto React
(\texttt{AuthContext}), que expone el usuario actual y la función de logout.

Las páginas implementadas son:
\begin{itemize}
  \item \textbf{/login}: formulario de login. Al autenticarse exitosamente, redirige a
        \texttt{/feed} (compradores) o a \texttt{/listings} (vendedores) según el rol devuelto
        por la API.
  \item \textbf{/register}: formulario de registro con selector de rol. Redirige al login tras
        el registro exitoso.
  \item \textbf{/feed}: feed de vehículos con swipe. Muestra una tarjeta a la vez con las fotos
        del vehículo y sus datos. El botón \emph{Like} llama a \texttt{POST /swipes} y muestra
        una notificación de match si la API indica que se creó uno nuevo. El botón \emph{Pass}
        también registra el swipe y avanza al siguiente vehículo.
  \item \textbf{/matches}: lista los matches del usuario. Cada tarjeta muestra el ID del match
        y la fecha de creación, con un enlace al chat correspondiente.
  \item \textbf{/chat/[id]}: chat de un match. Muestra la lista de mensajes y un formulario para
        enviar nuevos. Los mensajes se actualizan cada 3 segundos mediante polling HTTP. Se usa
        \texttt{useRef} para rastrear el conteo de mensajes sin agregar dependencias reactivas
        al intervalo, evitando reinicios innecesarios del timer.
  \item \textbf{/listings}: lista de mis publicaciones (solo vendedores). Permite navegar a la
        página de creación o de edición, y eliminar listings con confirmación.
  \item \textbf{/listings/new}: formulario de nueva publicación con subida de fotos
        (\texttt{PhotoUpload}), campos de marca, modelo, año, precio, tipo y descripción.
  \item \textbf{/listings/[id]/edit}: formulario de edición de un listing existente. Carga los
        datos actuales desde la API y permite modificarlos.
  \item \textbf{/profile}: vista del perfil del usuario autenticado con sus preferencias.
  \item \textbf{/profile/edit}: formulario para actualizar las preferencias de búsqueda del
        comprador (tipo de vehículo y rango de presupuesto).
\end{itemize}

El componente \texttt{PhotoUpload} permite seleccionar múltiples fotos, convertirlas a base64
mediante \texttt{FileReader} y previsualizarlas antes de publicar. Utiliza \texttt{Promise.all}
para procesar todos los archivos de forma concurrente y actualizar el estado en un único paso,
evitando el bug de closure que ocurre al leer archivos en un bucle con actualizaciones de estado
independientes.

\subsection{Infraestructura}

\begin{itemize}
  \item \textbf{PostgreSQL 16}: tablas \texttt{users}, \texttt{buyer\_profiles},
        \texttt{listings}, \texttt{listing\_photos}, \texttt{swipes}, \texttt{matches} y
        \texttt{messages}, con migraciones SQL versionadas en \texttt{apps/server/migrations/}.
  \item \textbf{Docker Compose}: orquesta los servicios \texttt{api} (Go, puerto 8080) y
        \texttt{db} (PostgreSQL, puerto 5432). El servicio \texttt{api} se construye desde el
        \texttt{Dockerfile} del servidor y depende de que \texttt{db} esté saludable.
\end{itemize}

% ===== 4. REQUERIMIENTOS Y REGLAS DE NEGOCIO =====
\section{Requerimientos y reglas de negocio}

\subsection{Requerimientos funcionales implementados}

\begin{center}
\begin{tabular}{cp{10.5cm}}
  \toprule
  \textbf{ID} & \textbf{Descripción} \\
  \midrule
  FR-01 & Registro de usuarios con selección de rol (comprador o vendedor). El email se
          normaliza a minúsculas y se valida su formato; la contraseña se hashea con bcrypt. \\[3pt]
  FR-02 & Autenticación mediante email y contraseña. Retorna un JWT (HS256) con \texttt{user\_id}
          y \texttt{role} en los claims, con 24 horas de expiración. \\[3pt]
  FR-03 & Feed personalizado para compradores. Los listings se filtran por estado
          \texttt{active}, por las preferencias del perfil del comprador (tipo y presupuesto) y
          excluyendo los ya swipeados y los publicados por el propio comprador. \\[3pt]
  FR-04 & Mecanismo de swipe (\emph{like}/\emph{pass}). Solo compradores pueden swipear listings
          activos que no sean propios. Cada swipe se persiste en la tabla
          \texttt{swipes}. \\[3pt]
  FR-12 & Gestión completa de listings por vendedores: crear (con fotos), editar campos
          parcialmente (PATCH), cambiar estado (\texttt{active}/\texttt{paused}/\texttt{sold})
          y eliminar. Solo el vendedor dueño puede modificar o eliminar su propio listing. \\[3pt]
  FR-19 & Creación automática de match al registrar un swipe de tipo \emph{like} sobre un
          listing activo. El match queda asociado al comprador y al listing. \\[3pt]
  FR-20 & Chat privado dentro de un match. Solo los participantes (comprador y vendedor del
          listing) pueden enviar y leer mensajes. \\
  \bottomrule
\end{tabular}
\end{center}

\subsection{Reglas de negocio}

Las siguientes reglas están implementadas y se verifican en la capa de dominio o de servicio,
nunca solo en el cliente:

\begin{enumerate}
  \item Solo los usuarios con rol \textbf{vendedor} (\texttt{seller}) pueden crear, editar y
        eliminar listings. \textit{Verificado en \texttt{ListingService.Create()} al consultar el
        rol del usuario antes de llamar a \texttt{NewListing()}.}

  \item Solo los usuarios con rol \textbf{comprador} (\texttt{buyer}) pueden consumir el feed y
        realizar swipes. \textit{Verificado en \texttt{FeedService.Build()} y en
        \texttt{MatchService.Swipe()} al comprobar \texttt{u.Role != domain.RoleBuyer}.}

  \item Un comprador no puede hacer swipe sobre un listing cuyo \texttt{seller\_id} coincida con
        su propio ID. \textit{Verificado en \texttt{MatchService.Swipe()} con la condición
        \texttt{listing.SellerID == buyerID}.}

  \item El match se crea automáticamente solo si el swipe es de tipo \emph{like} y el listing
        tiene estado \texttt{active}. Un listing en estado \texttt{paused} o \texttt{sold} no
        puede recibir swipes.

  \item Solo los participantes de un match pueden leer y enviar mensajes. \texttt{MatchService}
        verifica que el \texttt{userID} sea el \texttt{buyer\_id} del match o el
        \texttt{seller\_id} del listing asociado antes de permitir cualquier operación de chat.

  \item Un listing requiere al menos una foto para poder ser creado. \textit{Validado en
        \texttt{domain.NewListing()} con la condición \texttt{len(photoURLs) == 0}.}

  \item La contraseña debe tener un mínimo de 6 caracteres. \textit{Validado en
        \texttt{AuthService.Register()} antes de llamar al hasher.}

  \item El precio de un listing debe ser mayor a cero (\texttt{price > 0}). \textit{Validado en
        \texttt{domain.NewListing()} y en \texttt{domain.Listing.ApplyUpdates()}.}

  \item El año del vehículo, si se especifica, debe estar en el rango 1900--2100. \textit{Validado
        en \texttt{domain.NewListing()} y en \texttt{domain.Listing.ApplyUpdates()}.}

  \item La marca y el modelo son campos obligatorios y no pueden quedar vacíos tras aplicar
        \texttt{strings.TrimSpace()}. \textit{Validado en \texttt{domain.NewListing()} y en
        \texttt{domain.Listing.ApplyUpdates()}.}
\end{enumerate}

% ===== 5. PRINCIPIOS SOLID =====
\clearpage
\section{Principios SOLID}

\subsection{SRP --- Responsabilidad Única (\textit{Single Responsibility Principle})}

Cada componente tiene una única razón para cambiar. La separación en capas Handler / Service /
Repository es la materialización más directa de este principio, pero también se aplica dentro
de cada capa:

\begin{itemize}
  \item \textbf{AuthHandler}: su única responsabilidad es recibir peticiones HTTP de
        autenticación, validar el formato de los datos de entrada (\texttt{RegisterRequest},
        \texttt{LoginRequest}) y devolver las respuestas HTTP correctas. No contiene ninguna
        regla de negocio.
  \item \textbf{AuthService}: orquesta el proceso de registro (hashear contraseña, crear entidad
        User, persistir) y de login (buscar usuario, verificar contraseña, emitir JWT). No
        conoce ni HTTP ni SQL.
  \item \textbf{BcryptHasher}: única responsabilidad: aplicar el algoritmo bcrypt para hashear y
        verificar contraseñas. No sabe de usuarios ni de tokens.
  \item \textbf{JWTIssuer}: única responsabilidad: emitir tokens JWT con los claims
        \texttt{user\_id} y \texttt{role}, y validarlos al recibir peticiones.
  \item \textbf{FeedService}: única responsabilidad: construir el feed de un comprador concreto
        a partir de sus preferencias de perfil. Tiene un solo método público \texttt{Build()}.
  \item \textbf{postgresListingRepository}: única responsabilidad: ejecutar consultas SQL sobre
        las tablas \texttt{listings} y \texttt{listing\_photos}. Implementa la carga por lote
        de fotos para evitar el problema N+1, pero no contiene lógica de negocio.
\end{itemize}

\begin{center}
\begin{tikzpicture}[node distance=1.3cm, >=Stealth]
  \node[mybox, fill=green!15, minimum width=2.6cm] (h)
    {Handler\\\footnotesize HTTP / serialización};
  \node[mybox, fill=clssblue, minimum width=2.6cm, right=of h] (s)
    {Service\\\footnotesize Lógica de negocio};
  \node[mybox, fill=orange!20, minimum width=2.8cm, right=of s] (r)
    {Repository\\\footnotesize Persistencia SQL};
  \node[mybox, fill=gray!20, minimum width=2.2cm, right=of r] (db)
    {PostgreSQL\\\footnotesize Datos};
  \draw[flow] (h) -- (s);
  \draw[flow] (s) -- (r);
  \draw[flow] (r) -- (db);
\end{tikzpicture}
\end{center}

\subsection{OCP --- Abierto/Cerrado (\textit{Open/Closed Principle})}

El struct \texttt{ListingFilter} en \texttt{domain/listing.go} permite extender los criterios de
búsqueda sin modificar la firma de \texttt{ListingRepository.Search()} ni la lógica de
\texttt{FeedService}. Cualquier nuevo filtro (año mínimo, marca específica, radio geográfico)
se agrega como un nuevo campo puntero en el struct; si el campo es \texttt{nil}, el repositorio
simplemente no aplica esa condición. La función \texttt{Search()} nunca necesita cambiar su
signatura.

\begin{lstlisting}[style=gostyle]
// domain/listing.go
type ListingFilter struct {
    Status          *ListingStatus
    VehicleType     *string
    PriceMin        *int
    PriceMax        *int
    ExcludeSwipedBy *uuid.UUID  // excluye listings ya swipeados por el buyer
    ExcludeSellerID *uuid.UUID  // excluye listings del propio comprador
    Limit           int
    Offset          int
    // Nuevo filtro: agregar YearMin *int no modifica Search()
}
\end{lstlisting}

El repositorio construye la cláusula \texttt{WHERE} de forma dinámica inspeccionando qué campos
son no-\texttt{nil}, por lo que el OCP se mantiene tanto en la interfaz como en la
implementación.

\subsection{LSP --- Sustitución de Liskov (\textit{Liskov Substitution Principle})}

Toda implementación concreta de repositorio es sustituible por cualquier otra que cumpla la misma
interfaz sin que los servicios lo detecten. Por ejemplo, \texttt{postgresUserRepository} y un
hipotético \texttt{inMemoryUserRepository} (usado en tests unitarios) implementan la misma
interfaz \texttt{UserRepository}:

\begin{lstlisting}[style=gostyle]
// repository/user_repo.go
type UserRepository interface {
    Create(ctx context.Context, u *domain.User) error
    FindByEmail(ctx context.Context, email string) (*domain.User, error)
    FindByID(ctx context.Context, id uuid.UUID) (*domain.User, error)
}
\end{lstlisting}

\texttt{AuthService} recibe un \texttt{UserRepository} en su constructor y opera exactamente
igual independientemente de si la implementación concreta accede a PostgreSQL, a una base
en memoria o a un mock de test. El contrato de la interfaz (parámetros, valores de retorno,
semántica de errores como \texttt{ErrNotFound}) es idéntico en todas las implementaciones.

\subsection{ISP --- Segregación de Interfaces (\textit{Interface Segregation Principle})}

En lugar de una única interfaz \texttt{Repository} con todos los métodos del sistema, se definen
interfaces granulares por agregado. Cada servicio declara solo las dependencias que realmente
necesita:

\begin{center}
\begin{tabular}{lp{7.5cm}}
  \toprule
  \textbf{Interfaz} & \textbf{Métodos} \\
  \midrule
  \texttt{UserRepository}    & \texttt{Create}, \texttt{FindByEmail}, \texttt{FindByID} \\
  \texttt{ListingRepository} & \texttt{Create}, \texttt{FindByID}, \texttt{Search},
                               \texttt{ListBySeller}, \texttt{Update}, \texttt{Delete} \\
  \texttt{MatchRepository}   & \texttt{Create}, \texttt{FindByID}, \texttt{ListForUser} \\
  \texttt{MessageRepository} & \texttt{Create}, \texttt{ListByMatch} \\
  \texttt{ProfileRepository} & \texttt{Get}, \texttt{Upsert} \\
  \texttt{SwipeRepository}   & \texttt{Create} \\
  \bottomrule
\end{tabular}
\end{center}

Como consecuencia directa, \texttt{FeedService} solo recibe tres repositorios en su constructor
(\texttt{UserRepository}, \texttt{ProfileRepository} y \texttt{ListingRepository}), sin ninguna
dependencia sobre \texttt{MatchRepository}, \texttt{MessageRepository} ni
\texttt{SwipeRepository}. Si esas interfaces cambian, \texttt{FeedService} no se ve afectado.

\begin{lstlisting}[style=gostyle]
// service/feed_service.go
type FeedService struct {
    users    repository.UserRepository    // solo necesita FindByID
    profiles repository.ProfileRepository // solo necesita Get
    listings repository.ListingRepository // solo necesita Search
}
\end{lstlisting}

\subsection{DIP --- Inversión de Dependencias (\textit{Dependency Inversion Principle})}

Los módulos de alto nivel (servicios) no dependen de implementaciones concretas; dependen de
abstracciones (interfaces). Las implementaciones concretas tampoco dependen de los servicios;
ambas dependen de las interfaces definidas en el paquete \texttt{repository} y \texttt{auth}.

El \textbf{Composition Root} en \texttt{main.go} es el único lugar donde se instancian las
implementaciones concretas y se inyectan por constructor, manteniendo el resto del código libre
de acoplamientos:

\begin{lstlisting}[style=gostyle]
// cmd/api/main.go  (Composition Root)
hasher     := auth.NewBcryptHasher(bcrypt.DefaultCost)
jwtIssuer  := auth.NewJWTIssuer(cfg.JWTSecret, 24*time.Hour)
userRepo   := repository.NewPostgresUserRepository(conn)
// ...
authService := service.NewAuthService(userRepo, hasher, jwtIssuer)
// AuthService solo conoce las interfaces UserRepository y PasswordHasher
\end{lstlisting}

\begin{center}
\begin{tikzpicture}[node distance=1cm and 2.1cm, >=Stealth]
  \node[classbox] (as) {\texttt{AuthService}};

  \node[ifcbox, right=2.3cm of as, yshift=1cm, minimum width=3.8cm] (ph)
    {\texttt{PasswordHasher}\\{\scriptsize\textit{interface}}};
  \node[ifcbox, right=2.3cm of as, yshift=-1cm, minimum width=3.8cm] (ur)
    {\texttt{UserRepository}\\{\scriptsize\textit{interface}}};

  \node[classbox, right=2cm of ph, minimum width=3.4cm] (bh) {\texttt{BcryptHasher}};
  \node[classbox, right=2cm of ur, minimum width=3.4cm] (pg) {\texttt{postgresUserRepo}};

  \draw[uses] (as) -- (ph) node[midway,above,font=\scriptsize\sffamily]{usa};
  \draw[uses] (as) -- (ur);
  \draw[impl] (bh) -- (ph) node[midway,above,font=\scriptsize\sffamily]{implementa};
  \draw[impl] (pg) -- (ur) node[midway,above,font=\scriptsize\sffamily]{implementa};
\end{tikzpicture}
\end{center}

% ===== 6. PATRONES DE DISEÑO =====
\clearpage
\section{Patrones de diseño}

\subsection{Factory Method}

El patrón \emph{Factory Method} se aplica en las funciones constructoras del dominio y de la
capa de autenticación. Estas funciones actúan como fábricas que centralizan la validación de
invariantes de negocio: el llamador no puede construir un objeto en estado inválido porque la
única forma de obtenerlo es a través del constructor, que retorna un error descriptivo si alguna
invariante se viola.

\subsubsection{Instancias en el código}

\textbf{\texttt{domain.NewListing()}} --- valida que la marca no esté vacía, el modelo no esté
vacío, el precio sea positivo, el año (si se especifica) esté en el rango 1900--2100 y que se
incluya al menos una foto. Si todas las validaciones pasan, crea el \texttt{Listing} con estado
\texttt{active}, asigna un UUID nuevo y establece las marcas de tiempo:

\begin{lstlisting}[style=gostyle]
// domain/listing.go
func NewListing(sellerID uuid.UUID, brand, model string, year *int,
    price int, vehicleType, description string,
    photoURLs []string) (*Listing, error) {
    brand = strings.TrimSpace(brand)
    model = strings.TrimSpace(model)
    if brand == ""           { return nil, ErrInvalidBrand }
    if model == ""           { return nil, ErrInvalidModel }
    if price <= 0            { return nil, ErrInvalidPrice }
    if year != nil && (*year < 1900 || *year > 2100) {
        return nil, ErrInvalidYear
    }
    if len(photoURLs) == 0   { return nil, ErrPhotosRequired }
    // ... construye y retorna *Listing en estado "active"
}
\end{lstlisting}

\textbf{\texttt{domain.NewUser()}} --- normaliza el email a minúsculas y recorta espacios,
valida que contenga \texttt{@}, que el rol sea \texttt{buyer} o \texttt{seller} y que el hash de
contraseña no esté vacío.

\textbf{\texttt{auth.NewBcryptHasher(cost int)}} --- si el \texttt{cost} proporcionado está
fuera del rango válido de bcrypt (\texttt{MinCost}--\texttt{MaxCost}), lo reemplaza por
\texttt{bcrypt.DefaultCost}, garantizando un \texttt{BcryptHasher} siempre operable.

\textbf{\texttt{auth.NewJWTIssuer(secret string, ttl time.Duration)}} --- si el TTL es cero o
negativo, lo reemplaza por 24 horas como valor por defecto seguro.

\begin{center}
\begin{tikzpicture}[node distance=1.1cm and 2.4cm, >=Stealth]
  \node[mybox, fill=green!15, minimum width=3.5cm] (f)
    {\texttt{NewListing(\ldots)}\\\footnotesize Factory Method};
  \node[classbox, right=2.5cm of f, minimum width=3cm] (ok)
    {\texttt{*Listing}\\\footnotesize estado \texttt{active}, UUID nuevo};
  \node[mybox, fill=red!15, below=1.3cm of f, minimum width=3.5cm] (err)
    {\texttt{error}\\\footnotesize \texttt{ErrInvalidBrand}, \texttt{ErrInvalidPrice}, \ldots};
  \draw[flow] (f) -- (ok)  node[midway,above,font=\scriptsize\sffamily]{[válido]};
  \draw[flow] (f) -- (err) node[midway,right,font=\scriptsize\sffamily]{[invariante violada]};
\end{tikzpicture}
\end{center}

\subsection{Repository}

El patrón \emph{Repository} desacopla completamente la lógica de negocio del mecanismo de
persistencia. Los servicios nunca importan \texttt{database/sql} ni conocen el dialecto SQL; solo
operan con las interfaces de repositorio definidas en el paquete \texttt{repository}.

Cada agregado de dominio tiene su propia interfaz con los métodos estrictamente necesarios para
ese agregado. La implementación concreta (\texttt{postgres*Repository}) resuelve los detalles
de SQL, incluyendo optimizaciones como la carga por lote de fotos en \texttt{ListingRepository}
para evitar el problema N+1 (una consulta por todas las fotos de varios listings en lugar de una
consulta por cada listing):

\begin{lstlisting}[style=gostyle]
// repository/listing_repo.go
type ListingRepository interface {
    Create(ctx context.Context, l *domain.Listing) error
    FindByID(ctx context.Context, id uuid.UUID) (*domain.Listing, error)
    Search(ctx context.Context, f domain.ListingFilter) ([]*domain.Listing, error)
    ListBySeller(ctx context.Context, sellerID uuid.UUID) ([]*domain.Listing, error)
    Update(ctx context.Context, l *domain.Listing) error
    Delete(ctx context.Context, id uuid.UUID) error
}
\end{lstlisting}

\begin{center}
\begin{tikzpicture}[node distance=1.8cm and 2.2cm, >=Stealth]
  \node[ifcbox, minimum width=5cm, minimum height=1.6cm] (uri) {
    \texttt{ListingRepository}\\
    {\footnotesize\textit{interface}}\\
    {\scriptsize\texttt{Create / FindByID / Search}}\\
    {\scriptsize\texttt{ListBySeller / Update / Delete}}
  };
  \node[classbox, minimum width=5cm, minimum height=1.4cm, below=2cm of uri] (pg) {
    \texttt{postgresListingRepository}\\
    {\scriptsize\texttt{db: *sql.DB}}\\
    {\scriptsize carga fotos en lote (N+1 avoidance)}
  };
  \node[classbox, minimum width=3cm, left=2.2cm of uri] (svc)
    {\texttt{ListingService}\\{\footnotesize y \texttt{FeedService}}};

  \draw[impl]  (pg)  -- (uri) node[midway,right,font=\scriptsize\sffamily]{implementa};
  \draw[uses]  (svc) -- (uri) node[midway,above,font=\scriptsize\sffamily]{depende de};
\end{tikzpicture}
\end{center}

\subsection{Strategy}

El patrón \emph{Strategy} se aplica en el mecanismo de hashing de contraseñas. La interfaz
\texttt{PasswordHasher} define el contrato con dos métodos: \texttt{Hash()} para hashear una
contraseña en texto plano y \texttt{Verify()} para compararla contra un hash almacenado.

\texttt{AuthService} recibe un \texttt{PasswordHasher} en su constructor y nunca instancia
ninguna implementación concreta. Esto significa que el algoritmo de hashing puede cambiarse
(por ejemplo, de bcrypt a argon2 o scrypt) simplemente implementando la interfaz y pasando la
nueva implementación al constructor de \texttt{AuthService}, sin modificar una sola línea de
la lógica de autenticación:

\begin{lstlisting}[style=gostyle]
// auth/hasher.go
// PasswordHasher define la estrategia para hashear contraseñas.
// Patrón Strategy: AuthService depende de esta interface, no de bcrypt.
type PasswordHasher interface {
    Hash(plain string) (string, error)
    Verify(plain, hashed string) error
}

type BcryptHasher struct { cost int }

func (h *BcryptHasher) Hash(plain string) (string, error) {
    b, err := bcrypt.GenerateFromPassword([]byte(plain), h.cost)
    return string(b), err
}

func (h *BcryptHasher) Verify(plain, hashed string) error {
    return bcrypt.CompareHashAndPassword([]byte(hashed), []byte(plain))
}
\end{lstlisting}

\begin{center}
\begin{tikzpicture}[node distance=1.8cm and 2.2cm, >=Stealth]
  \node[ifcbox, minimum width=5cm, minimum height=1.8cm] (ph) {
    \texttt{PasswordHasher}\\
    {\footnotesize\textit{interface (Strategy)}}\\
    {\scriptsize\texttt{Hash(plain) (string, error)}}\\
    {\scriptsize\texttt{Verify(plain, hashed) error}}
  };
  \node[classbox, minimum width=5cm, minimum height=1.4cm, below=2cm of ph] (bh) {
    \texttt{BcryptHasher}\\
    {\scriptsize\textit{estrategia concreta}}\\
    {\scriptsize\texttt{cost: int} --- algoritmo bcrypt}
  };
  \node[classbox, minimum width=3cm, left=2.2cm of ph] (as)
    {\texttt{AuthService}\\{\footnotesize cliente de la estrategia}};

  \draw[impl] (bh) -- (ph) node[midway,right,font=\scriptsize\sffamily]{implementa};
  \draw[uses] (as) -- (ph) node[midway,above,font=\scriptsize\sffamily]{usa};
\end{tikzpicture}
\end{center}

% ===== 7. APIs =====
\section{APIs}

El backend expone una API REST bajo el prefijo \texttt{/api/v1} usando el framework Gin 1.12.
La documentación interactiva completa está disponible en
\texttt{http://localhost:8080/swagger/index.html}, generada con \texttt{swaggo/gin-swagger}.

\subsection{Autenticación de requests protegidos}

Los endpoints protegidos requieren el header:
\begin{lstlisting}[style=gostyle, numbers=none]
Authorization: Bearer <jwt-token>
\end{lstlisting}

El JWT contiene los claims \texttt{user\_id} (UUID), \texttt{role} (\texttt{buyer} o
\texttt{seller}), \texttt{iat} (emitido en) y \texttt{exp} (expira en, 24 horas tras la emisión).
El middleware \texttt{JWTAuth} valida la firma con HS256 y rechaza tokens expirados o
mal formados con \texttt{401 Unauthorized}.

\subsection{Tabla de endpoints}

\begin{longtable}{lll}
  \toprule
  \textbf{Método} & \textbf{Ruta} & \textbf{Descripción} \\
  \midrule
  \endhead
  \multicolumn{3}{l}{\small\textit{Públicos (sin JWT)}} \\[3pt]
  GET    & \texttt{/api/v1/health}               & Estado del servidor (\texttt{200 OK}) \\
  POST   & \texttt{/api/v1/auth/register}        & Registro de nuevo usuario \\
  POST   & \texttt{/api/v1/auth/login}           & Login; retorna JWT y datos del usuario \\
  GET    & \texttt{/api/v1/listings}             & Búsqueda pública de listings \\
  GET    & \texttt{/api/v1/listings/:id}         & Detalle de un listing \\
  \midrule
  \multicolumn{3}{l}{\small\textit{Protegidos (requieren JWT)}} \\[3pt]
  GET    & \texttt{/api/v1/profile/me}           & Obtener perfil del usuario autenticado \\
  PUT    & \texttt{/api/v1/profile/me}           & Actualizar preferencias del perfil \\
  POST   & \texttt{/api/v1/listings}             & Crear listing (solo \texttt{seller}) \\
  GET    & \texttt{/api/v1/listings/me}          & Listar mis listings \\
  PATCH  & \texttt{/api/v1/listings/:id}         & Editar listing (solo su dueño) \\
  DELETE & \texttt{/api/v1/listings/:id}         & Eliminar listing (solo su dueño) \\
  GET    & \texttt{/api/v1/feed}                 & Feed personalizado (solo \texttt{buyer}) \\
  POST   & \texttt{/api/v1/swipes}               & Registrar swipe like/pass (solo \texttt{buyer}) \\
  GET    & \texttt{/api/v1/matches}              & Listar mis matches \\
  GET    & \texttt{/api/v1/matches/:id/messages} & Listar mensajes de un match \\
  POST   & \texttt{/api/v1/matches/:id/messages} & Enviar mensaje en un match \\
  \bottomrule
\end{longtable}

\subsection{Esquemas de request y response}

\subsubsection{Registro e inicio de sesión}

\begin{lstlisting}[style=gostyle, numbers=none]
// POST /api/v1/auth/register
{ "email": "usuario@ejemplo.cl", "password": "secreto", "role": "buyer" }
// Response 201: { "id": "uuid", "email": "...", "role": "buyer" }

// POST /api/v1/auth/login
{ "email": "usuario@ejemplo.cl", "password": "secreto" }
// Response 200:
{ "token": "eyJ...", "user": { "id": "uuid", "email": "...", "role": "buyer" } }
\end{lstlisting}

\subsubsection{Listings}

Todos los endpoints de listado (\texttt{/listings}, \texttt{/listings/me}, \texttt{/feed})
retornan la misma estructura paginada con el campo \texttt{items} y el conteo total:

\begin{lstlisting}[style=gostyle, numbers=none]
// GET /api/v1/listings  |  GET /api/v1/listings/me  |  GET /api/v1/feed
{
  "items": [
    {
      "id": "uuid", "seller_id": "uuid",
      "brand": "Toyota", "model": "Corolla",
      "year": 2020, "price": 12000000,
      "vehicle_type": "sedan", "description": "...",
      "status": "active",
      "photos": [{ "id": "uuid", "url": "...", "position": 0 }],
      "created_at": "2025-01-01T00:00:00Z",
      "updated_at": "2025-01-01T00:00:00Z"
    }
  ],
  "count": 1
}
\end{lstlisting}

\texttt{GET /api/v1/feed} acepta los parámetros de consulta \texttt{limit} (máx.\,100,
por defecto 20) y \texttt{offset} para paginación.

\subsubsection{Perfil y matches}

\begin{lstlisting}[style=gostyle, numbers=none]
// GET /api/v1/profile/me
{ "user_id": "uuid", "vehicle_type": "sedan",
  "budget_min": 5000000, "budget_max": 15000000,
  "updated_at": "2025-01-01T00:00:00Z" }

// GET /api/v1/matches  (estructura paginada igual a listings)
{
  "items": [
    { "id": "uuid", "buyer_id": "uuid",
      "listing_id": "uuid", "created_at": "2025-01-01T00:00:00Z" }
  ],
  "count": 1
}

// GET /api/v1/matches/:id/messages
{
  "items": [
    { "id": "uuid", "match_id": "uuid",
      "sender_id": "uuid", "body": "Hola, ¿sigue disponible?",
      "created_at": "2025-01-01T00:00:00Z" }
  ],
  "count": 1
}
\end{lstlisting}

% ===== 8. BACKLOG Y PARTICIPACIÓN =====
\section{Backlog y participación del equipo}

\subsection{Historial de commits}

\begin{longtable}{lp{7.5cm}l}
  \toprule
  \textbf{Commit} & \textbf{Descripción} & \textbf{Autor} \\
  \midrule
  \endhead
  \texttt{7006743}
    & Inicializar repositorio en rama main; estructura base del monorepo
    & B. Aliaga \\[3pt]
  \texttt{0cfdb3c}
    & Backend primera entrega: dominio (User, BuyerProfile), Auth service y handler,
      Profile service y handler, repositorios Postgres, migraciones iniciales, Docker
      Compose, Swagger
    & B. Aliaga \\[3pt]
  \texttt{279a2b8}
    & Backend: módulo completo de listings --- dominio Listing/Photo, ListingRepository
      con carga por lote, ListingService con validación de rol, handler CRUD (FR-12)
    & B. Aliaga \\[3pt]
  \texttt{31801e4}
    & Backend: feed con filtro por preferencias (FR-03), swipe con match automático
      (FR-04, FR-19), chat dentro de match (FR-20), MatchService con verificación de
      participante
    & B. Aliaga \\[3pt]
  \texttt{aee1550}
    & Frontend: inicialización del proyecto cliente con Next.js 16, TypeScript,
      Tailwind CSS y estructura de carpetas
    & L. Salazar / J. Carrera \\[3pt]
  \texttt{ec120d5}
    & Frontend: página de perfil con links a gestión de vehículos y edición de
      preferencias; integración con AuthContext
    & L. Salazar / J. Carrera \\[3pt]
  \texttt{dc085d5}
    & Frontend: \texttt{package.json} con dependencias y devDependencies del proyecto
    & B. De La Fuente \\[3pt]
  \texttt{86e5fd0}
    & Frontend: flujo completo de registro de usuario y hooks de autenticación
      (\texttt{AuthContext}, \texttt{saveAuthData}, \texttt{ProtectedRoute})
    & B. De La Fuente \\[3pt]
  \texttt{69b86b4}
    & Frontend: páginas de feed con swipe, matches, chat con polling, edición de
      listings y rutas protegidas por rol
    & L. Salazar / J. Carrera \\[3pt]
  \texttt{0b97781}
    & Frontend: mejoras en la creación de listings (PhotoUpload con Promise.all),
      correcciones de contratos API y manejo de errores sin \texttt{alert()}
    & L. Salazar / J. Carrera \\
  \bottomrule
\end{longtable}

\subsection{Distribución del trabajo}

\begin{center}
\begin{tabular}{lp{9.5cm}}
  \toprule
  \textbf{Integrante} & \textbf{Responsabilidades principales} \\
  \midrule
  Benjamín Aliaga Mardones
    & Diseño y desarrollo completo del backend (Go): arquitectura por capas, entidades de
      dominio con constructores validadores, interfaces de repositorio e implementaciones
      PostgreSQL, lógica de servicios (auth, listings, feed, match, chat, perfil),
      handlers HTTP con Gin, migraciones SQL, Docker Compose, configuración de Swagger
      y CORS. \\[5pt]
  Lizardo Salazar \& Juan Carrera
    & Desarrollo colaborativo del frontend (Next.js): inicialización del proyecto,
      integración de autenticación JWT con contexto React, implementación de todas las
      páginas (feed con swipe, matches, chat con polling, listings CRUD, perfil y
      preferencias), componente PhotoUpload, rutas protegidas por rol y ajuste de los
      contratos de API con el backend. \\[5pt]
  Benjamín De La Fuente
    & Configuración del entorno frontend (\texttt{package.json}), implementación del
      flujo de registro de usuario y desarrollo de los hooks y componentes de
      autenticación (\texttt{AuthContext}, \texttt{ProtectedRoute},
      \texttt{saveAuthData}). \\
  \bottomrule
\end{tabular}
\end{center}

\end{document}
